\documentclass[12pt]{article}
\usepackage[utf8]{inputenc}
\usepackage[T1]{fontenc}
\usepackage{hyperref}
\title{Sanapan\'a Language Resources for K-12 Curriculum Development}
\author{Compiled by Claude Sonnet 5.5 (claude-sonnet-5-5) for Andr\'as Kornai}
\date{}
\begin{document}
\maketitle

\section{Introduction}
The code \texttt{sap} (Glottocode sana1281) is listed as retired. The language it denotes, Sanapan\'a, is now covered by ISO 639-3 \texttt{spn}. Sanapan\'a is an Enlhet-Enenlhet (Mascoyan) language of the Paraguayan Chaco. Its speakers have also been called Sanapana, Quilyacmoc and Kelya'mok, and they call themselves Enenlhet or nenlhet in some sources. Estimates of speakers are about 980 (2007) and about 914 who speak it at home, with a 2002 census count of about 2,270 ethnic Sanapan\'a. Endangered Languages reports that in La Esperanza and Anaconda nearly the whole community still speaks it and children still acquire it at home. The searches found no school materials in Sanapan\'a, and the list below is correspondingly short.

\section{Primary Educational Resources}
I found no primers, textbooks or school curricula in Sanapan\'a. The closest items are the documentation and description projects listed below, which could supply source material for teaching materials.
\begin{itemize}
\item A documentation of historical narratives among the Sanapan\'a (Culture in Crisis; ELDP-funded project led by Jens Van Gysel, fieldwork in La Esperanza): \url{https://cultureincrisis.org/projects/a-documentation-of-historical-narratives-amongst-the-sanapaná-enlhet-enenlhet-of-the-paraguayan-chaco}
\end{itemize}

\section{Linguistic and Reference Materials}
\begin{itemize}
\item Van Gysel, Jens L. (2017), \emph{Temporal predicative particles in Sanapan\'a and the Enlhet-Enenlhet language family}, Leiden MA thesis; bibliographic record: \url{https://glottolog.org/resource/reference/id/687722}
\item Van Gysel's University of New Mexico dissertation, ``A Usage-Based Account of Ongoing Structural Changes in Sanapan\'a: Grammar Sketch and Case Studies'' (2024); I saw it described in search results but did not see a direct URL.
\item Silva Gomes, Antonio Almir (2013), \emph{Sanapan\'a uma l\'ingua Maskoy: aspectos gramaticais}, PhD thesis, UNICAMP (Portuguese); I saw it described in search results but did not see a direct URL.
\item English Wikipedia overview: \url{https://en.wikipedia.org/wiki/Sanapan%C3%A1_language}
\item Intercontinental Dictionary Series word lists for related Enlhet and Angait\'e varieties are mentioned on Wikipedia, but I did not confirm their contents.
\end{itemize}

\section{Religious and Community Resources}
No Sanapan\'a scripture or community-organisation resources were found in the searches.

\section{Digital and Audio Resources}
\begin{itemize}
\item Endangered Languages Project entry for Sanapan\'a: \url{https://endangeredlanguages.com/lang/2752}
\item Van Gysel's narratives project (above) deposits recordings in the Endangered Languages Archive (ELAR). I could not confirm from search results that the materials are already available.
\end{itemize}

\section{Curriculum Development Context}
Sanapan\'a is still transmitted to children in at least two communities, but there is a shift toward Guaran\'i and Spanish. No Sanapan\'a-medium schooling materials were found. A curriculum would have to start from the new documentation (narratives, grammar sketch) and be developed with the communities. Contacting the documentation team and the communities is the sensible first step.

\end{document}